\documentclass{article}
\usepackage{hyperref}
\usepackage[margin=0.5in]{geometry}
\usepackage{amsmath}
\usepackage{graphicx}
\usepackage[bottom]{footmisc}
\usepackage{lastpage}
\usepackage{fancyhdr}
\usepackage{float}
\usepackage{tabularx}

\graphicspath{ {./assets/} }

\pagestyle{fancy}
\fancyhf{} % clear existing header/footer entries
% Place Page X of Y on the right-hand
% side of the footer
\fancyfoot[R]{Yang \thepage \hspace{1pt} of \pageref{LastPage}}
\renewcommand{\headrulewidth}{0pt}

\title{PHYSIOL 578: HW 2.4\footnotetext{U-M GPT was used to help draft and refine the answers.}}
\author{Jack Yang \\ \href{mailto:yhyunjoo@umich.edu}{yhyunjoo@umich.edu}}

\begin{document}
\maketitle
\thispagestyle{fancy}

\section*{
    Q1.
    You performed mass spectrometry on immunopurified mTOR kinase from HEK293T cells and identified a novel protein (TIP1; TOR Interacting Protein 1).
    Sequence analysis revealed that TIP1 is a unique protein with a molecular mass of 35 kDa, and no similar proteins to TIP1 have been identified in the human genome.
    You successfully generated an antibody that specifically recognizes TIP1 (35 kD protein) in western blotting using HEK293T cell lysates.
    However, the TIP1 antibody does not work for immunoprecipitation. \\
    \\
    How would you determine which mTOR complex contains TIP1? \\
    \\
    \underline{Provide one key experiment and describe how you would interpret the results.} \\
    \\
    (For your experiments, assume you have antibodies (for western blotting and immunoprecipitation experiments) for the components of the mTOR complexes. (2 points)
}

Perform co-immunoprecipitation under lysis conditions that preserve mTOR complexes, such as CHAPS. Immunoprecipitate Raptor to isolate mTORC1 and, in parallel, Rictor to isolate mTORC2. Include a nonspecific IgG immunoprecipitation as a negative control. Then immunoblot each precipitate using the TIP1 antibody. Blotting for mTOR, Raptor, and Rictor can confirm that the appropriate complexes were successfully isolated.

If TIP1 co-immunoprecipitates with Raptor but not Rictor, TIP1 associates with mTORC1.
If TIP1 co-immunoprecipitates with Rictor but not Raptor, TIP1 associates with mTORC2.
If TIP1 co-immunoprecipitates with both, it may associate with both complexes.
If TIP1 is detected in neither, its interaction may be weak or disrupted during lysis, or it may associate with a different mTOR-containing complex.

The absence of TIP1 from the IgG control would support the specificity of the observed association.

\section*{
    Q2.
    You performed mass spectrometry analyses of the immunopurified GATOR2 from HEK293T cells treated with a glucose-free medium and identified a novel GATOR2-interacting protein (GG2IP; Glucose-sensitive GATOR2 Interacting Protein).
    You failed to detect GG2IP peptides in mass spectrometry analyses when you prepared the sample under glucose (500 mg/dl)-replete conditions.
    Knockdown of endogenous GG2IP using several shRNAs consistently enhanced S6K1 and 4EBP1 phosphorylation in HEK293T cells under glucose-starved conditions, but had no effect under glucose-replete conditions.
    In contrast, Akt phosphorylation was enhanced in HEK293T cells lacking GG2IP expression under glucose-starved conditions. \\
    \\
    Describe the roles of GG2IP in the regulation of mTORC1 and mTORC2 activity.
}

GG2IP appears to be a glucose-dependent negative regulator of both mTORC1 and mTORC2 signaling. Under glucose-replete conditions, GG2IP knockdown has no detectable effect. Under glucose-starved conditions, GG2IP associates with GATOR2 and normally suppresses mTORC1 activity, as shown by the increased phosphorylation of S6K1 and 4E-BP1 following GG2IP knockdown. One possible mechanism is that GG2IP inhibits GATOR2, thereby allowing GATOR1 to suppress mTORC1. GG2IP also normally suppresses mTORC2-dependent AKT phosphorylation during glucose starvation, because loss of GG2IP increases AKT phosphorylation.


\end{document}